\subsection{External source and frozen chronological design}\label{sec:external-design}
We test whether an executable acquisition policy reaches a justified model
comparison with fewer requested records on a source not used in the earlier
energy experiments. Low Carbon London (LCL) provides half-hour energy records
from UK Power Networks through the London Datastore \cite{lcl}. We scan the
complete official archive and take the 20 lowest household IDs with at least
90 complete training days among its lowest 100 IDs. Selection uses training
availability alone, without forecast or acquisition outcomes. The sampling
frame is deterministic and is not a representative population sample.
One selected household has no eligible test days and is retained in the
attrition inventory rather than replaced. All subsequently eligible days
and all six pairs of the four candidate models enter the analysis.

An hourly reference is the sum of the two finite, nonnegative source values
labeled at its :00 and :30 positions. Both must be present and unique;
duplicates, nulls, negative entries and incomplete hours are excluded, with
no interpolation, replacement by zero or outlier clipping. A complete day
requires all 24 references; feature availability imposes an additional,
predeclared restriction. The source has no actual-versus-estimated quality
flag, so this screen does not identify every possible provider-side
estimation practice. The provider's physical timezone and interval-endpoint
convention are not established by the inspected metadata. We therefore map
its labels onto a declared 24-hour benchmark clock; UTC-shaped timestamps in
the implementation are experiment coordinates, not verified physical UTC.
The two 2013 clock-change dates are conservatively excluded. Training is
January--September 2013, validation October--December 2013, and testing
January--February 2014. Source records, selection and code hashes were
frozen before the corresponding evaluation; this is a local protocol freeze,
not a public preregistration.

The original external-source plan used REFIT raw power records \cite{refitraw}. Under its
frozen strict reference-quality rule, all 20 houses yielded no complete
days, so no REFIT forecast or recovery outcome was evaluated. We retained
that failed feasibility inventory and selected the LCL fallback before
model fitting or acquisition outcomes. No reference-quality rule was relaxed
to obtain a favorable result.

\subsection{Causal forecasts and two permitted archives}\label{sec:external-archive}
The candidates are persistence, standardized Ridge regression, histogram
gradient boosting and the fixed multilayer perceptron. Learned-model
specifications, iteration allowances and random states are unchanged from
the historical study in Appendix~\ref{app:historical-design}; each selected
household receives its own training-only fit. The new causal feature vector
contains corresponding-hour values from days $D-2,D-3,D-8$, the mean of day
$D-2$, and sine/cosine hour, weekday and month terms. Forecasts for day $D$
are assigned issue time noon on $D-1$. Feature-source records end no later
than the start of $D-1$; learned forecasts also require that the last
training target precede their issue time. Persistence copies the
corresponding hour of $D-2$. Forecast generation does not require its own
future verification target: available persistence rows are retained even
when learned features or future reference targets are incomplete.

For each target day the early audit cutoff is the start of $D+1$ and the
later cutoff is noon on $D+1$, immediately after that issue. Every saved
prediction permitted by the cutoff is available, including already-issued
predictions for future targets. At the early cutoff, predictions involving
day $D$ as a lag have not yet been issued. At the later cutoff, the saved
persistence forecasts for $D+2$ reveal every target in $D$. Reconstruction
uses those forecast values and source-row traces, not hidden reference
values. The historical targets are used afterward to check the copies.
Because no predictor issued by the early cutoff uses day $D$ as an input,
inverting these learned forecasts cannot add constraints on its missing
targets either, conditional on this fixed-model benchmark schedule.
Both cutoffs, target availability and record access are constructed
benchmark rules. They do not measure production arrival times or establish
a physical twelve-hour service-level benefit.

Each day has four exact six-hour totals constructed from its eligible
references. A single contiguous six-hour gap has a uniformly selected start
from 0 through 18, using three fixed SHA-256-derived streams per day. Masks
apply to the audit inputs, not to past model training or issued features.
Thus the benchmark represents an incomplete audit export while the forecast
service may later access the underlying readings; it does not simulate a
persistent physical-meter failure shared by both systems.
All policies receive the same initially surviving readings, totals and
permitted archive. Direct copies, singleton residuals and zero residual
totals are incorporated for free before any request. These controlled totals
are not evidence that independent registers survived real LCL outages.

\subsection{Executable stopping policies and equal information}\label{sec:external-policies}
Five policies operate sequentially on each fixed model pair: exact
deterministic width minimization, greedy width reduction, a Hara-inspired
proposal, uniform sampling and the randomized robust method with its
uniform-barycenter secondary optimum. At an unresolved stage each requests
at most two currently unknown coordinates, reveals the entire batch, applies
free aggregate deductions and recomputes the comparison interval. It stops
at a strict certificate, a numerical tie, or six cumulative requested
coordinates. Every submitted coordinate is charged, including both members
of a batch that already suffices for certification; there are no refunds or
duplicate requests. Already certified comparisons use no requests. Each
policy excludes already determined and contrast-irrelevant coordinates
from its acquisition candidates under the same rule. Each
coordinate is an hourly target record composed of two source rows, not a
measured meter transaction. Requests remain pair-specific; they are not
shared across the six model pairs.

The Hara-inspired proposal adapts the squared-loss comparison weight
$|a_i-b_i|\mathbb{E}|a_i+b_i-2Y_i|$ from \cite{hara}. A fixed Ridge mean and
validation-residual Gaussian scale specify a proposal distribution;
conditioning independent Gaussian errors on each residual total gives its
mean and marginal variance. The absolute-value expectation is evaluated
analytically. Nonnegative support is enforced by the certificate, not by
this Gaussian proposal. We sample without replacement with these weights,
using uniform weights if all vanish. This is an explicitly described
regression adaptation, not a reproduction of the published classification
experiment or its importance-weighted estimator. No test labels calibrate
the surrogate. Validation acquisition replay uses the full-validation
residual scale and is therefore a retrospective calibration diagnostic;
the test run is chronological because the scale is frozen beforehand.

The exact batch method minimizes worst remaining width at each stage; it
does not solve a minimum-cost sequential stopping problem. Likewise the
randomized policy optimizes worst expected width, not expected requests to
certification. The three fixed streams define three actual acquisition
paths, not an exact expectation over all randomized trajectories. Every
policy sees the same available observations and uses the same enclosure and
strict-sign tests. Numerical ties remain uncertified and are reported rather
than resolved through hidden reference values.
If the deterministic routine returns an empty optimal batch despite positive
width, we select the lexicographically first feasible full batch with the
same primary optimum. This outcome-independent tie rule prevents a stalled
acquisition sequence; it does not improve the claimed minimax value.

\subsection{Predeclared decision-cost outcomes}\label{sec:external-metric}
The primary curve is the fraction certified by cumulative integer budget
$B=0,\ldots,6$ among initially uncertified comparisons, including numerical
ties. Fractions are first computed within each household and then averaged
equally across households with a nonzero denominator. We also report pooled
counts, the all-case curve, learned-model-only pairs, each household, each
model pair and each test month. Mean request costs include failures and
early stops. The descriptive budget $B_{80}$ is the smallest integer whose
primary coverage reaches 80\%; if none does, it is not reached. Curves count
only actually executed batches and do not interpolate fractional requests.

Before outcomes we declared a practical planning criterion: at least a
15\% reduction in $B_{80}$ relative to the better of greedy and Hara, without
lower coverage at the six-request cap. Uniform and randomized results are
reported as well, so an uncompetitive baseline cannot be hidden by that
comparison. This criterion is a study-planning threshold, not a claim of
statistical significance. Whole-household,
pair and learned-only diagnostics reveal whether an aggregate result depends
on a few cases. No household replacement, hyperparameter tuning or acquisition
retuning is allowed after the test outcomes are opened.
